\section{引言}

模仿学习能够从人类示范中获得复杂机器人操作策略，但其性能高度依赖训练数据的规模与覆盖范围。对于双臂灵巧操作，这一数据瓶颈更加突出：操作者必须同时控制两个机械臂和夹爪或灵巧手，示范采集负担高；与此同时，插入、穿线、交接和装配等任务又对接触状态及左右臂同步关系十分敏感。MimicGen~\cite{mandlekar2023mimicgen} 和 DexMimicGen~\cite{jiang2024dexmimicgen} 表明，可以从少量源示范出发，在仿真中自动生成大量机器人轨迹。然而，双臂数据生成仍需要回答一个关键问题：\emph{哪些相对时序变化属于有效多样性，哪些变化会破坏任务所依赖的物理协同？}

这一问题不能仅由运动学可达性回答。即使左右臂的每个末端目标都可达，组合轨迹仍可能因接触时刻错位、支撑手提前离开、操作手过早施力或物体失稳而失败。换言之，\textbf{运动学上可达不代表物理上可执行}。现有的空间变换、噪声注入和统一时间缩放，往往同时改变轨迹几何、动作幅度或全局速度，因此很难将失败归因于双臂相对时序本身。

本文研究少量成功示范下的双臂协同时序增强问题。我们的基本思路是构造受控反事实：保持某一手臂原始动作的顺序和内容不变，只在指定阶段插入等待，使该手臂的后缀相对另一手臂发生延迟。由于没有删除动作，该干预避免了“时序变化”和“累计动作缺失”之间的主要混杂。需要强调的是，等待期间环境状态仍会自然演化，因此本文所称“反事实”是可执行的干预式反事实，而不是所有变量均被固定的严格因果反事实。我们随后在物理仿真中执行轨迹，以稳定成功、任务距离和接触统计评价反事实结果，进而得到示范的协同敏感图谱。

与只保留成功合成轨迹的数据生成范式不同，本文同时利用两类反事实样本：接近原示范且仍成功的样本用于扩展策略的时序容忍范围；仅因相对时序变化而失败的困难负样本用于学习协同边界。由此，研究问题从“能否产生更多轨迹”转化为“能否产生对双臂协同学习真正有信息量的轨迹”。

本文的预期贡献如下：
\begin{enumerate}
    \item 提出一种保序、动作完整的双臂时序反事实干预，降低动作删除对相对时序分析的混杂；
    \item 提出基于仿真执行、稳定成功和接触几何的协同敏感性度量，用于定位关键手臂、关键阶段与延迟区间；
    \item 构建面向少样本模仿学习的反事实数据选择与训练方案，并在多个数据规模和任务上验证其相对普通噪声增强、统一时间缩放和未筛选延迟增强的优势；
    \item \todo{若完成实机实验：验证仿真中识别的协同敏感性及策略增益能否迁移到真实双臂平台。}
\end{enumerate}

\paragraph{论文组织。}
第2节回顾双臂模仿学习、自动数据生成和时序建模；第3节介绍问题定义、反事实生成与协同感知训练；第4节给出实验问题、基线和评价指标；第5节总结全文。
